\section*{第 \ref{ch:g3:shapes} 章——图形与直角}

\begin{solution}{exo:g3:shapes:1}
典型的\omterm{def:g3:shapes:rightangle}{直角}有：黑板的角、一张纸的角、
门的角。非\omterm{def:g3:shapes:rightangle}{直角}：一把剪刀张开时的角、
1 点时钟的指针所成的角。每个角都要用
三角尺（\cref{met:g3:shapes:check}）检查——绝不能只凭眼睛。
\end{solution}

\begin{solution}{exo:g3:shapes:2}
作法：先画一条长 $6$ cm 的边，并在两端用三角尺画\omterm{def:g3:shapes:rightangle}{直角}，
再画长 $4$ cm 的边，并封闭图形。标记：$4$ 个小方框
画在四个角上，每条 $6$ cm 的边画一道短线，
每条 $4$ cm 的边画两道短线。
\end{solution}

\begin{solution}{exo:g3:shapes:3}
正方形有 $4$ 个\omterm{def:g3:shapes:rightangle}{直角}和 $4$ 条相等的边，图中都已标出。
\end{solution}

\begin{solution}{exo:g3:shapes:4}
方格作图：方格保证了\omterm{def:g3:shapes:rightangle}{直角}，
数方格可以保证边长正确。
\end{solution}

\begin{solution}{exo:g3:shapes:5}
“每个正方形都是长方形”：\emph{正确}——它有 $4$ 个\omterm{def:g3:shapes:rightangle}{直角}，
这正是长方形所要求的全部。

“每个长方形都是正方形”：\emph{错误}——一个 $6 \times 4$
长方形的边长不全相等。

“三角形可以有两个\omterm{def:g3:shapes:rightangle}{直角}”：\emph{错误}——从两个\omterm{def:g3:shapes:rightangle}{直角}出发
画出的两条边，都与同一条底边所在的直线垂直，
因此永远不会相交，无法封闭成三角形。
\end{solution}

\begin{solution}{exo:g3:shapes:6}
$OA = 4$ cm（即\omterm{def:g3:shapes:shapes}{半径}）。$OB = 4$ cm 也是，因为圆上\emph{每一点}
到圆心的距离都相同。
\end{solution}

\begin{solution}{exo:g3:shapes:7}
这个长方形被分成三个三角形：两个\omterm{def:g3:shapes:rightangle}{直角}三角形
（$AMD$ 和 $MBC$，分别以 $A$ 和 $B$ 为\omterm{def:g3:shapes:rightangle}{直角}）以及中间的
三角形 $MCD$。
\end{solution}

\begin{solution}{exo:g3:shapes:8}
至少含有一个\omterm{def:g3:shapes:rightangle}{直角}的图形有：正方形、长方形、\omterm{def:g3:shapes:rightangle}{直角}
三角形，以及字母 L（它的内、外拐角）。圆
没有任何角。
\end{solution}

\begin{solution}{exo:g3:shapes:9}
每条边都沿“向右 $2$ 格、向上 $1$ 格”（或旋转后的
“向右 $1$ 格、向下 $2$ 格”等走法）穿过，每条边都是
一个 $2 \times 1$ 方格块的对角线，因此四条边长度相同。
\end{solution}

\begin{solution}{exo:g3:shapes:10}
一个正确的作法示例：“1. 画长方形 $ABCD$，使 $AB = 3$ cm
（底边），$BC = 2$ cm。2. 在边 $[DC]$ 上方画
正方形 $DCEF$，边长为 $3$ cm，画在长方形外侧。”任何清楚、
编号齐全、步骤完整的作法都是正确的——
逐字照做即可检验。
\end{solution}

\begin{solution}{exo:g3:shapes:11}
$1 \times 1$ 的正方形：$9$ 个。$2 \times 2$ 的正方形：$4$ 个。
$3 \times 3$：$1$ 个。共有 $9 + 4 + 1 = 14$ 个正方形。
\end{solution}
